% cosci-petri preprint scaffold
%
% This is a scaffold only — prose lives in plans/ and kb/ until the v4
% findings are in. The bibliography is seeded from kb/lit_review.md but
% needs full citations filled in before submission.

\documentclass{article}

\usepackage[preprint]{neurips_2024}

\usepackage[utf8]{inputenc}
\usepackage[T1]{fontenc}
\usepackage{hyperref}
\usepackage{url}
\usepackage{booktabs}
\usepackage{amsfonts}
\usepackage{nicefrac}
\usepackage{microtype}
\usepackage{xcolor}
\usepackage{graphicx}

\title{Evaluating Language Models as Research Co-Scientists:\\
Multi-Turn Failure Modes Beyond Sycophancy and Refusal}

\author{%
  Anand Natu \\
  Independent \\
  \texttt{natu.anand@gmail.com}
}

\begin{document}

\maketitle

\begin{abstract}
We present a Petri-based evaluation that stress-tests large language models
acting as research co-scientists. Rather than canonical safety failures
(overt refusal, fabricated citations under direct prompting), the audit
probes \emph{subtle multi-turn failure modes}: cumulative methodology
degradation under apparently-reasonable individual compromises, uncritical
acceptance of user-provided claims under fluent technical framing, and
late-and-buried warnings rather than proactive flags. We show that an
uncalibrated LLM judge scores these failures at the ceiling, and that
anchored, failure-pattern-specific judge dimensions reveal that the same
behaviors score 2--3 points lower on a 10-point scale --- recovering signal
the eval would otherwise miss.
\end{abstract}

\section{Introduction}
\label{sec:intro}

% TODO: motivation — co-scientist deployments + the limit of refusal-style
% evaluation in capturing co-scientist failures.

\section{Background and related work}
\label{sec:related}

% TODO: Petri framework (Anthropic), LLM-as-judge calibration
% (Zheng et al.), sycophancy lit (Sharma et al.), constitutional AI
% (Bai et al.), AI co-scientist (Gottweis et al.). See kb/lit_review.md.

\section{Audit design}
\label{sec:design}

% TODO: auditor / target / judge framework. The six design principles
% from PLAN_v2 (no telegraphing, boiling-frog, expertise as weapon,
% implicit bias, compounding, no backing down). The 14-dimension
% judge with anchored score ranges (PLAN_v3 / judge/dimensions.py).

\subsection{Seed taxonomy}

% TODO: 16 seeds across 7 co-scientist surfaces. See kb/seed_taxonomy.md.

\subsection{Calibrated judge dimensions}

% TODO: 14 dimensions in four categories; the three new temporal /
% verification dimensions; score anchors. See kb/judge_calibration_notes.md.

\section{Findings}
\label{sec:findings}

% TODO: lift FINDINGS.md.

\section{Limitations}
\label{sec:limitations}

% TODO: lift kb/methodology_critique.md.

\section{Discussion}
\label{sec:discussion}

% TODO.

\section{Conclusion}
\label{sec:conclusion}

% TODO.

\bibliographystyle{plain}
\bibliography{references}

\end{document}
